\documentclass[12pt,a4paper]{article}
\usepackage[utf8]{inputenc}
\usepackage{amsmath, amssymb}
\usepackage{geometry}
\usepackage{setspace}
\usepackage{graphicx}
\usepackage{listings}
\usepackage{xcolor}

\geometry{margin=2.5cm}
\setstretch{1.2}

\definecolor{codebg}{gray}{0.97}
\lstset{
    language=Python,
    backgroundcolor=\color{codebg},
    basicstyle=\ttfamily\small,
    keywordstyle=\color{blue},
    commentstyle=\color{green!50!black},
    numbers=left,
    numberstyle=\tiny,
    frame=single,
    breaklines=true
}

\title{Taller de Álgebra Lineal}
\author{José Luciano Mejía Arias}
\date{}

\makeatletter
\renewcommand{\maketitle}{
  \begin{flushright}
    \includegraphics[width=3cm]{../../logo_unal.png}
  \end{flushright}
  \begin{center}
    {\LARGE \textbf{\@title}}\\[0.3cm]
    {\large \@author}
  \end{center}
}
\makeatother

\begin{document}
\maketitle

\section*{Instrucciones}
Resuelva cada ejercicio mostrando el procedimiento completo. Es importante justificar cada paso y escribir la respuesta final con claridad.

\section*{Enunciado}

\section{Sistemas de ecuaciones y forma matricial}
Escriba cada sistema en forma matricial $AX=b$ y luego indique la matriz aumentada correspondiente.

\begin{enumerate}
    \item
    \[
    \begin{cases}
    2x + 3y = 7 \\
    x - y = 1
    \end{cases}
    \]

    \item
    \[
    \begin{cases}
    3x - 2y = 4 \\
    x + y = 5
    \end{cases}
    \]

    \item
    \[
    \begin{cases}
    x + y + z = 6 \\
    2x - y + z = 3 \\
    x + 2y - z = 4
    \end{cases}
    \]
\end{enumerate}

\section{Eliminación de Gauss}
Resuelva los siguientes sistemas utilizando el método de Gauss.

\begin{enumerate}
    \item
    \[
    \begin{cases}
    x + 2y = 5 \\
    3x - y = 1
    \end{cases}
    \]

    \item
    \[
    \begin{cases}
    2x + y = 8 \\
    x - 3y = -7
    \end{cases}
    \]

    \item
    \[
    \begin{cases}
    x + y + z = 6 \\
    2x - y + 3z = 14 \\
    3x + 2y - z = 2
    \end{cases}
    \]
\end{enumerate}

\section{Determinantes}
Calcule el determinante de cada matriz.

\begin{enumerate}
    \item
    \[
    A = \begin{pmatrix}
    2 & 3 \\
    5 & 7
    \end{pmatrix}
    \]

    \item
    \[
    B = \begin{pmatrix}
    1 & 2 \\
    3 & 4
    \end{pmatrix}
    \]

    \item
    \[
    C = \begin{pmatrix}
    3 & 1 & 2 \\
    2 & 0 & 1 \\
    1 & 4 & 3
    \end{pmatrix}
    \]
\end{enumerate}

\section*{Ejercicio básico en Python}
Escriba un programa en Python que pida el nombre, apellido y edad de una persona. Si la persona es mayor de edad, debe mostrar un mensaje de bienvenida. Si no lo es, debe mostrar un mensaje de bloqueo.

\begin{lstlisting}[language=Python]
# Ejercicio en Python
# Pedir nombre, apellido y edad
# Si edad >= 18:
#     mostrar "Bienvenido"
# else:
#     mostrar "No puedes entrar"
\end{lstlisting}

\end{document}
